\section{Data collection}

The substrates that were chosen to characterize the market of the Cremona province for demonstrative purposes are reported in Table \ref{tab:items} (Appendix), along with their legislative cluster, expected/deterministic value of the technical characteristics, transportation fixed tariff, purchase price and BMP$_{\infty}$.
The clusters are products (P), industrial by-products (B), animal slurries (S) and second-harvest crops (P2).

The yearly availability of animal and agro-industrial residues was taken from the ISTAT (italian national statistics institute) database. Data for animal residues (slurries and dungs) were available for the 2018-2020 time period. Similarly, data for agro-industrial residues (silages etc...) were available for the 2016-2020 time period. 

A BMP$_{\infty}$ database was built concatenating some literature values with the ones from two main sources: PoliMi lab-tests and AgriPower (A2A S.p.A Group). In the latter, each data point has also its TS and VS value recorder. Some data points presented also a TKN value. The other item characteristics were taken from other literature sources. The time-dependent BMP curve was created assuming different kinetic rates (d$^{-1}$) and asymptotic value BMP$_{\infty}$ of each item.

Substrate prices were retrieved from multiple sources, primarily from: i) the \emph{Granaria Milano} website (Chamber of Commerce data), ii) the \emph{TESEO} web application, iii) old AgriPower plant business plans that were made available to the author. In source 'i' and 'ii', the data points for the primary product (grain) of maize, sorghum and barley cultivation was present for each month of 2009, 2014, 2019-2023. Nevertheless, the price of silages is in general equal to their production cost, so that source 'iii' was used to set the expected value and 'max/mean'/'min/mean' ratios were taken from the higher abundance of data of the other sources, to retrieve realistic distributions.
Only the distribution of 'wheat by-products' was taken as the original one from market data.
For other substrates, in face to the lack of data, only minimum, maximum and mean values were retrieved from source 'iii'.
For tomato peels, no data was found, so that symmetric distrubition with minimum and maximum values equal to 1 and 1.5\% of the tomato supermarket retailing was hypothesized. 

The criteria for the selection of a subset of substrates were i) availability $\geq$ 100 tonFM/year and ii) at least two BMP$_{\infty}$ data available. In the end, \emph{n}=10 substrates were selected.

Since no data were available for the spatial distribution of the sellers within the Cremona, distance values from each of the \emph{m}=20 sellers to the plant was hypothesized; in addition, 14 of them were considered to be 'farmers', whereas 2 and 3 were considered  to be tomato sauce (production of 'tomato peels') and wheat ('wheat by-products') industrial producers. One player is represented by the plant itself, which produces 'cow slurry','cow dung' and 'maize silage'. Also, the latter was considered for a scenario in which, at the end of the previous 'purchase horizon', some stocked substrate is remained unused. 
A uniform distribution was considered to randomly allocate the overall available quantities of slurries and energy crops among the players. 

The collected data were thus exploited for the solution of the deterministic linear problem (D-LP) and for the identification of reasonable distributions from which to generate the data points used for later design of uncertainty sets.

The reactor volume was considered equal to the one of a standard 1 MWe plant. 
